\documentclass[12pt]{article}
\usepackage{amsmath}
\usepackage{amssymb}
\usepackage{graphicx}
\usepackage{geometry}
\geometry{a4paper, margin=1in}

\usepackage{fontspec}
\setmainfont{Times New Roman}

\title{\textbf{Intermediate Microeconomics Practice 1}}
\author{Mochiao Chen}
\date{\today}

\begin{document}

\maketitle

\noindent \textbf{Instructions:}
\begin{itemize}
    \item Answer all questions in English.
    \item Write your calculation process clearly for computational problems.
    \item For graphing questions, clearly label all axes, curves, and equilibrium points.
\end{itemize}

\section*{Part I: True or False (20 points)}
\textit{Indicate whether the following statements are True (T) or False (F). Briefly explain your reasoning for any False statements.}

\begin{enumerate}
    \item \textbf{[ \quad ]} Walras' Law states that if $n-1$ markets in an exchange economy are in equilibrium, the $n$-th market must also be in equilibrium.
    \item \textbf{[ \quad ]} In a First-Degree (Perfect) Price Discrimination scenario, the monopolist extracts all consumer surplus, resulting in zero deadweight loss.
    \item \textbf{[ \quad ]} A pure public good is characterized by being rival in consumption and excludable.
    \item \textbf{[ \quad ]} In the Prisoner's Dilemma game, the Nash Equilibrium is always Pareto efficient.
    \item \textbf{[ \quad ]} Adverse selection usually occurs \textit{before} the transaction takes place due to hidden information, whereas moral hazard occurs \textit{after} the transaction due to hidden action.
    \item \textbf{[ \quad ]} Under Third-Degree Price Discrimination, a monopolist will charge a higher price in the market with more elastic demand (more negative price elasticity).
    \item \textbf{[ \quad ]} According to the Coase Theorem, if property rights are well-defined and transaction costs are zero, the initial allocation of property rights determines whether the efficient outcome can be achieved.
    \item \textbf{[ \quad ]} The "Core" of an exchange economy is the set of all Pareto-optimal allocations that are welfare-improving for both consumers relative to their initial endowments.
    \item \textbf{[ \quad ]} In a natural monopoly, setting the price equal to Marginal Cost ($P = MC$) allows the firm to break even (zero economic profit).
    \item \textbf{[ \quad ]} In a mixed-strategy Nash Equilibrium, a player chooses probabilities such that the opponent is indifferent between their pure strategies.
\end{enumerate}

\section*{Part II: Multiple Choice (20 points)}
\textit{Select the best answer for each question.}

\begin{enumerate}
    \item A monopolist faces a demand curve with constant elasticity $\varepsilon = -3$. If the marginal cost of production is constant at \$20, what is the profit-maximizing price?
    \begin{itemize}
        \item[A.] \$20
        \item[B.] \$30
        \item[C.] \$40
        \item[D.] \$60
    \end{itemize}

    \item Which of the following best describes the "Tragedy of the Commons"?
    \begin{itemize}
        \item[A.] A failure of public goods provision due to free-riding.
        \item[B.] Over-consumption of a non-excludable but rival resource.
        \item[C.] Under-production of a good with positive externalities.
        \item[D.] A coordination failure in a Stag Hunt game.
    \end{itemize}

    \item In an Edgeworth Box, the Contract Curve consists of:
    \begin{itemize}
        \item[A.] All allocations where one consumer is better off than the other.
        \item[B.] All allocations where the marginal rate of substitution (MRS) is equal for both consumers.
        \item[C.] All allocations inside the "lens" formed by the initial indifference curves.
        \item[D.] Only the competitive equilibrium points.
    \end{itemize}

    \item Consider the "Market for Lemons" (Adverse Selection). If buyers cannot distinguish between high-quality and low-quality cars, and only know the average value:
    \begin{itemize}
        \item[A.] Only high-quality cars will be traded.
        \item[B.] Both types will be traded at the average price.
        \item[C.] High-quality car sellers may exit the market, leading to a market collapse.
        \item[D.] Signaling has no effect on this market.
    \end{itemize}

    \item The efficiency condition for the provision of a public good with two consumers (A and B) is given by:
    \begin{itemize}
        \item[A.] $MRS_A = MRS_B = MRT$
        \item[B.] $MRS_A + MRS_B = MC$ (or $MRT$)
        \item[C.] $MU_A = MU_B$
        \item[D.] $MRS_A / P_A = MRS_B / P_B$
    \end{itemize}

    \item In a Sequential Game where Firm A moves first (Stackelberg leader) and Firm B moves second (follower):
    \begin{itemize}
        \item[A.] Firm A usually achieves a higher profit than in a simultaneous Cournot game.
        \item[B.] Firm B has the "first-mover advantage."
        \item[C.] The outcome is identical to the Cournot Nash Equilibrium.
        \item[D.] Both firms produce the monopoly quantity.
    \end{itemize}

    \item If a monopolist implements a Two-Part Tariff ($T = F + p \cdot q$) to maximize total profit and extract all surplus from identical consumers, it should set the per-unit price $p$ equal to:
    \begin{itemize}
        \item[A.] The monopoly price.
        \item[B.] Zero.
        \item[C.] Marginal Cost.
        \item[D.] Average Total Cost.
    \end{itemize}

    \item High-ability workers getting a college degree to distinguish themselves from low-ability workers, even if the degree adds no productivity, is an example of:
    \begin{itemize}
        \item[A.] Moral Hazard
        \item[B.] Screening
        \item[C.] Signaling
        \item[D.] Bundling
    \end{itemize}

    \item In the presence of a negative production externality (e.g., pollution), the private market equilibrium quantity is usually \_\_\_\_\_\_ than the socially optimal quantity.
    \begin{itemize}
        \item[A.] Lower
        \item[B.] Higher
        \item[C.] The same
        \item[D.] Unrelated
    \end{itemize}

    \item A situation where a consumer's utility depends on the consumption of others, or a production function depends on the activities of others, without market compensation, is defined as:
    \begin{itemize}
        \item[A.] An Externality
        \item[B.] A Public Good
        \item[C.] Imperfect Competition
        \item[D.] Asymmetric Information
    \end{itemize}
\end{enumerate}

\section*{Part III: Term Definitions (15 points)}
\textit{Briefly explain the following terms in the context of microeconomics (2-3 sentences each).}

\begin{enumerate}
    \item \textbf{Pareto Efficiency}
    \item \textbf{Nash Equilibrium}
    \item \textbf{Moral Hazard}
    \item \textbf{Third-Degree Price Discrimination}
    \item \textbf{Free-Rider Problem}
\end{enumerate}

\section*{Part IV: Calculations (30 points)}
\textit{Show all steps of your work.}

\subsection*{Question 1: Monopoly and Price Discrimination (10 points)}
A monopolist sells a product in two separate markets, Market 1 and Market 2. The demand curves are:
\[ P_1 = 100 - Q_1 \]
\[ P_2 = 80 - 2Q_2 \]
The monopolist's total cost function is $TC(Q) = 20Q$, where $Q = Q_1 + Q_2$.
\begin{enumerate}
    \item[(a)] If the monopolist can practice third-degree price discrimination, calculate the profit-maximizing quantity ($Q_1, Q_2$) and price ($P_1, P_2$) for each market.
    \item[(b)] Calculate the total profit under price discrimination.
    \item[(c)] Which market has a more elastic demand at the equilibrium point? Verify using the markup rule or elasticity formula.
\end{enumerate}

\subsection*{Question 2: Game Theory and Mixed Strategies (10 points)}
Consider the following payoff matrix for a game between Player A and Player B:

\begin{center}
\begin{tabular}{c|c|c}
 & \textbf{B: Left} & \textbf{B: Right} \\ \hline
\textbf{A: Up} & 2, 3 & 0, 1 \\ \hline
\textbf{A: Down} & 1, 0 & 3, 2 \\
\end{tabular}
\end{center}

\begin{enumerate}
    \item[(a)] Does this game have any pure-strategy Nash Equilibria? If yes, identify them. If no, explain why.
    \item[(b)] Find the mixed-strategy Nash Equilibrium. Let $p$ be the probability Player A plays Up, and $q$ be the probability Player B plays Left. Calculate the equilibrium values for $p$ and $q$.
\end{enumerate}

\subsection*{Question 3: Externalities (10 points)}
A steel factory (Firm S) is located upstream from a fishery (Firm F).
The steel firm's profit function is $\pi_S = 12s - s^2 - (x-4)^2$, where $s$ is steel output and $x$ is pollution.
The fishery's profit function is $\pi_F = 10f - f^2 - xf$, where $f$ is fish output.
The pollution $x$ negatively affects the fishery.

\begin{enumerate}
    \item[(a)] If the firms operate independently (steel firm ignores the external cost), what level of pollution $x$ will the steel firm choose?
    \item[(b)] If the two firms merge and maximize joint profit, what is the socially optimal level of pollution $x^*$?
    \item[(c)] Comparing (a) and (b), does the independent market produce too much or too little pollution?
\end{enumerate}

\section*{Part V: Graphing and Analysis (15 points)}

\textbf{Topic: General Equilibrium and the Edgeworth Box}

Consider an exchange economy with two consumers, A and B, and two goods, 1 and 2.
Consumer A has an initial endowment $\omega_A = (8, 2)$ and Consumer B has an initial endowment $\omega_B = (2, 8)$.
Their preferences are standard (convex indifference curves).

\begin{enumerate}
    \item[(a)] Draw an \textbf{Edgeworth Box} diagram. Clearly label the following:
    \begin{itemize}
        \item The dimensions of the box (total quantity of Good 1 and Good 2).
        \item The endowment point $W$.
        \item An indifference curve for A and an indifference curve for B passing through the endowment point.
    \end{itemize}
    \item[(b)] On your diagram, shade the area representing the \textbf{Lens of Trade} (the set of allocations that make both consumers better off).
    \item[(c)] Illustrate and label the \textbf{Contract Curve}. Explain what condition holds true at every point on the Contract Curve.
    \item[(d)] Identify the \textbf{Core} on the Contract Curve.
    \item[(e)] Draw a budget line and a potential competitive equilibrium point. Explain how the price ratio ($P_1/P_2$) relates to the indifference curves at this equilibrium.
\end{enumerate}

\end{document}